\section{Conclusão}

A avaliação primária por validação cruzada agrupada mostrou AUC ROC de
$0{,}9709\pm0{,}0059$ e AUC PR de $0{,}9607\pm0{,}0121$ para a ConvNeXt-Tiny na
triagem binária de células cervicais em citologia convencional. Esses valores
indicam boa capacidade de ranqueamento, mas a decisão final depende do limiar:
na partição retida, o regime de alta sensibilidade reduziu os falsos negativos
de 60 para 30 em relação ao limiar 0,50, ao custo de elevar os falsos positivos
de 64 para 146. Essa troca é compatível com triagem apenas quando há revisão
humana posterior.

A análise por categoria Bethesda mostrou que a simplificação binária não resolve
completamente a dificuldade das classes limítrofes. ASC-US foi a principal fonte de erro
observada, com 81,7\% de acerto entre 115 células, enquanto HSIL e SCC foram
mais facilmente recuperadas. A comparação exploratória com ResNet50 também
mostrou que a vantagem da ConvNeXt-Tiny está principalmente no ranqueamento, não
em todas as métricas dependentes de limiar.

As principais limitações são a ausência de identificadores de paciente na CRIC,
que não permite garantir separação clínica completa, a comparação com ResNet50
restrita à partição retida e a ausência de validação externa. Trabalhos futuros
devem repetir baselines nos cinco folds agrupados, validar o método em outras
bases, incorporar explicabilidade e avançar da célula recortada para localização
automática e agregação de evidências no nível da lâmina. Até essas etapas, o
método permanece um protótipo de pesquisa para priorização, sem finalidade
diagnóstica.
